% System Design

\section{System Design}
\label{sec:design}

This section presents the design of \textsc{WasmSentinel}, focusing on its three key innovations: the WasmML detection engine, the hybrid detection architecture, and the zero-copy data channel.

\subsection{Architecture Overview}
\label{sec:design-overview}

\begin{figure*}[t]
\centering
\begin{tikzpicture}[
    box/.style={rectangle, draw, minimum height=0.8cm, align=center, font=\footnotesize},
    ebpf/.style={box, fill=blue!15, minimum width=1.8cm},
    host/.style={box, fill=green!15, minimum width=2cm},
    wasm/.style={box, fill=orange!15, minimum width=2cm},
    arrow/.style={->, >=stealth, thick}
]
    % Kernel space
    \node[ebpf] (syscall) at (0,0) {Syscall\\Tracing};
    \node[ebpf] (sched) at (2,0) {Process\\Lifecycle};
    \node[ebpf] (net) at (4,0) {Network\\Events};
    \node[ebpf] (file) at (6,0) {File\\Access};
    \node[ebpf] (ring) at (3,-1.2) {Ring Buffer (BPF Map)};

    % User space
    \node[host] (reader) at (0,-3) {Event\\Reader};
    \node[host] (batcher) at (2.5,-3) {Batch\\Aggregator};
    \node[host] (engine) at (5,-3) {Wasm\\Engine};

    % Wasm modules
    \node[wasm] (stide) at (8,-2.5) {STIDE\\Detector};
    \node[wasm] (mlp) at (8,-3.5) {MLP\\Classifier};
    \node[wasm] (hybrid) at (10,-3) {Hybrid\\Scheduler};

    % Arrows
    \draw[arrow] (syscall) -- (ring);
    \draw[arrow] (sched) -- (ring);
    \draw[arrow] (net) -- (ring);
    \draw[arrow] (file) -- (ring);
    \draw[arrow] (ring) -- (reader);
    \draw[arrow] (reader) -- (batcher);
    \draw[arrow] (batcher) -- (engine);
    \draw[arrow] (engine) -- (stide);
    \draw[arrow] (engine) -- (mlp);
    \draw[arrow] (stide) -- (hybrid);
    \draw[arrow] (mlp) -- (hybrid);

    % Labels
    \node[above=0.5cm of net, font=\small\bfseries] {Kernel Space};
    \node[above=0.3cm of batcher, font=\small\bfseries] {User Space (Host)};
    \node[above=0.3cm of hybrid, font=\small\bfseries] {Wasm Sandbox};

    % Boundary lines
    \draw[dashed, thick, gray] (-1.5,-0.6) -- (11,-0.6);
    \draw[dashed, thick, gray] (6.5,-4.2) -- (6.5,-1.8);
\end{tikzpicture}
\caption{\textsc{WasmSentinel} architecture. eBPF programs collect events in kernel space and emit them to a ring buffer. The host program reads events and dispatches them to Wasm-sandboxed detection modules.}
\label{fig:architecture}
\end{figure*}

Figure~\ref{fig:architecture} illustrates the overall architecture of \textsc{WasmSentinel}.
The system comprises three layers:

\textbf{Kernel Layer.}
eBPF programs attach to kernel tracepoints to capture security-relevant events including syscall invocations, process lifecycle events, network connections, and file accesses.
Events are emitted to a shared ring buffer for efficient transfer to user space.

\textbf{Host Layer.}
A Go-based host program reads events from the ring buffer, optionally batches them for efficiency, and dispatches them to the Wasm detection engine.
The host also manages Wasm module lifecycle, resource limits, and alert handling.

\textbf{Detection Layer.}
Detection logic executes within Wasm sandboxes, providing strong isolation guarantees.
Our hybrid architecture combines multiple detection strategies, each implemented as a Rust module compiled to Wasm.

\subsection{WasmML: ML in Wasm Sandbox}
\label{sec:design-wasmml}

The core innovation of \textsc{WasmSentinel} is \textit{WasmML}---the ability to run machine learning inference inside a Wasm sandbox while maintaining microsecond-level latency.

\subsubsection{Design Challenges}

Running ML in Wasm presents several challenges:

\begin{enumerate}
    \item \textbf{No Native ML Libraries:} Standard frameworks (TensorFlow, PyTorch) are not available in Wasm
    \item \textbf{Limited Instruction Set:} Wasm lacks SIMD and other accelerator instructions (though SIMD is an emerging extension)
    \item \textbf{Memory Constraints:} Wasm linear memory must be pre-allocated
    \item \textbf{No Dynamic Loading:} Models must be embedded or loaded through host functions
\end{enumerate}

\subsubsection{Our Approach: Pure-Rust MLP}

Rather than attempting to port heavyweight ML frameworks, we implement a lightweight multi-layer perceptron (MLP) directly in Rust:

\begin{lstlisting}[style=rust, caption={Core MLP inference loop in Rust.}]
pub fn forward(&self, input: &[f32]) -> f32 {
    let mut x = input.to_vec();
    for (i, layer) in self.layers.iter().enumerate() {
        x = layer.weights.iter()
            .map(|row| {
                row.iter().zip(x.iter())
                   .map(|(w, x)| w * x)
                   .sum::<f32>()
            })
            .zip(layer.biases.iter())
            .map(|(sum, b)| {
                if i < self.layers.len() - 1 {
                    relu(sum + b)  // Hidden layers
                } else {
                    sigmoid(sum + b)  // Output
                }
            })
            .collect();
    }
    x[0]
}
\end{lstlisting}

This approach offers several advantages:
\begin{itemize}
    \item \textbf{Zero Dependencies:} No external ML libraries required
    \item \textbf{Predictable Performance:} No garbage collection or JIT compilation
    \item \textbf{Small Binary:} Complete detector compiles to 180KB Wasm
    \item \textbf{Portable Weights:} Model parameters stored as JSON, loaded at initialization
\end{itemize}

\subsubsection{Model Architecture}

Through systematic architecture search (Section~\ref{sec:eval-architecture}), we found that a 4-layer MLP achieves the best accuracy-latency trade-off:

\begin{table}[t]
\centering
\caption{MLP architecture for \textsc{WasmSentinel}.}
\label{tab:mlp-arch}
\begin{tabular}{lccc}
\toprule
\textbf{Layer} & \textbf{Input} & \textbf{Output} & \textbf{Activation} \\
\midrule
Input & 64 & 128 & ReLU \\
Hidden 1 & 128 & 64 & ReLU \\
Hidden 2 & 64 & 32 & ReLU \\
Hidden 3 & 32 & 16 & ReLU \\
Output & 16 & 1 & Sigmoid \\
\bottomrule
\end{tabular}
\end{table}

The 64-dimensional input feature vector encodes:
\begin{itemize}
    \item Syscall category (one-hot encoded)
    \item Dangerous syscall indicators (ptrace, execve, etc.)
    \item User/privilege level
    \item Process context features
    \item N-gram sequence statistics
\end{itemize}

Total model parameters: 19,201 (approximately 75KB when stored as float32).

\subsection{Hybrid Detection Architecture}
\label{sec:design-hybrid}

\begin{figure}[t]
\centering
\begin{tikzpicture}[
    box/.style={rectangle, draw, minimum height=0.7cm, align=center, font=\footnotesize},
    decision/.style={diamond, draw, aspect=2.5, font=\tiny, inner sep=1pt},
    arrow/.style={->, >=stealth}
]
    \node[box, fill=blue!10] (event) at (0,0) {Event};
    \node[box, fill=green!10] (stide) at (2,0) {STIDE};
    \node[decision, fill=yellow!10] (check) at (4.2,0) {Anomaly?};
    \node[box, fill=orange!10] (mlp) at (6.4,0.8) {MLP};
    \node[box, fill=green!10] (normal) at (6.4,-0.8) {Normal};
    \node[decision, fill=yellow!10] (mlpcheck) at (8.6,0.8) {Threat?};
    \node[box, fill=red!10] (alert) at (10.5,0.8) {Alert};
    \node[box, fill=green!10] (benign) at (10.5,0) {Benign};

    \draw[arrow] (event) -- (stide);
    \draw[arrow] (stide) -- (check);
    \draw[arrow] (check) -- node[above, font=\tiny] {Yes} (mlp);
    \draw[arrow] (check) -- node[right, font=\tiny] {No} (normal);
    \draw[arrow] (mlp) -- (mlpcheck);
    \draw[arrow] (mlpcheck) -- node[above, font=\tiny] {Yes} (alert);
    \draw[arrow] (mlpcheck) -- node[right, font=\tiny] {No} (benign);
\end{tikzpicture}
\caption{Hybrid detection flow: STIDE filters normal events, suspicious events escalate to MLP.}
\label{fig:hybrid-flow}
\end{figure}

Pure ML-based detection incurs inference cost for every event.
To reduce overhead, we employ a \textit{hybrid detection architecture} that combines STIDE-based fast-path filtering with MLP-based deep analysis.

\subsubsection{STIDE Fast Path}

STIDE (Sequence Time-Delay Embedding)~\cite{stide} is a lightweight anomaly detection algorithm based on n-gram analysis of syscall sequences.
During training, STIDE builds a database of ``normal'' n-grams observed in benign behavior.
At runtime, incoming syscall sequences are checked against this database; sequences containing novel n-grams are flagged as anomalous.

Our implementation uses:
\begin{itemize}
    \item N-gram size: 6 (empirically optimal for syscall sequences)
    \item Hash function: FNV-1a for fast, collision-resistant hashing
    \item Database size: Up to 100,000 unique n-grams
\end{itemize}

STIDE overhead is approximately 200ns per event, dominated by hash computation.

\subsubsection{Escalation Policy}

Events proceed through the detection pipeline as follows:

\begin{enumerate}
    \item \textbf{STIDE Check:} Compute anomaly score based on novel n-gram ratio
    \item \textbf{Fast Path:} If score $< \theta_1$ (default: 0.3), classify as benign
    \item \textbf{Escalation:} If score $\geq \theta_1$, extract features and invoke MLP
    \item \textbf{Deep Analysis:} MLP returns threat probability
    \item \textbf{Classification:} If probability $> \theta_2$ (default: 0.7), generate alert
\end{enumerate}

This strategy achieves the best of both worlds:
\begin{itemize}
    \item \textbf{Low overhead:} 80--90\% of benign events are filtered by STIDE
    \item \textbf{High accuracy:} Complex attacks are analyzed by MLP
    \item \textbf{Tunable trade-off:} Thresholds $\theta_1$ and $\theta_2$ control precision-recall balance
\end{itemize}

\subsection{Zero-Copy Data Channel}
\label{sec:design-zerocopy}

Efficient data transfer from kernel to Wasm is critical for overall performance.
Traditional approaches incur multiple memory copies and serialization overhead.

\subsubsection{Traditional Approach}

A naive implementation involves:
\begin{enumerate}
    \item eBPF writes event to ring buffer
    \item Host copies event to Go heap
    \item Go serializes event to JSON
    \item Host copies JSON to Wasm linear memory
    \item Wasm parses JSON to extract fields
\end{enumerate}

This incurs 4 memory copies and expensive JSON parsing, adding 5--50$\mu$s per event depending on event size.

\subsubsection{Our Approach: Direct Memory Mapping}

\textsc{WasmSentinel} reduces this to a single copy:
\begin{enumerate}
    \item eBPF writes event to ring buffer in a fixed binary format
    \item Host maps ring buffer region directly into Wasm linear memory (read-only)
    \item Wasm reads event fields at known offsets
\end{enumerate}

\begin{lstlisting}[style=rust, caption={Zero-copy event access in Wasm.}]
// Event structure matches kernel layout
#[repr(C)]
struct RawEvent {
    timestamp: u64,
    pid: u32,
    ppid: u32,
    uid: u32,
    event_type: u32,
    syscall_nr: u32,
    // ... additional fields
}

fn read_event(ptr: *const u8) -> &RawEvent {
    unsafe { &*(ptr as *const RawEvent) }
}
\end{lstlisting}

This approach reduces per-event overhead to under 500ns.

\subsubsection{Safety Considerations}

Direct memory mapping raises safety concerns that we address:
\begin{itemize}
    \item \textbf{Read-only access:} Wasm cannot modify the ring buffer
    \item \textbf{Bounds checking:} Event access is validated against buffer boundaries
    \item \textbf{Format validation:} Magic numbers and checksums detect corruption
    \item \textbf{Double-fetch protection:} Events are copied to stack before multi-field access
\end{itemize}
